\section{Torque surfaces and their classification}

The linear EESM torque is
\begin{equation}
 \Te=\kt\iq(\Lexc\iexc+\dL\id),\qquad
 \kt=\frac32p,\quad\dL=\Ld-\Lq.\label{eq:eesm-torque}
\end{equation}
Introduce the flux-producing coordinate and an independent complementary
coordinate
\begin{equation}
 \sflux=\dL\id+\Lexc\iexc,\qquad
 r=\Lexc\id-\dL\iexc.
\end{equation}

These coordinates can be discovered rather than guessed.  Torque already
contains the combination $\dL i_d+\Lexc i_e$, so naming that combination
$s_\psi$ isolates the direction that is effective for torque.  In the
$(i_d,i_e)$ plane its coefficient vector is $[\dL,\Lexc]\trans$.  A natural
complement is the perpendicular vector $[\Lexc,-\dL]\trans$, because
\begin{equation}
 [\dL,\Lexc]\begin{bmatrix}\Lexc\\-\dL\end{bmatrix}=0.
\end{equation}
This produces $r=\Lexc i_d-\dL i_e$.  In matrix form,
\begin{equation}
 \begin{bmatrix}s_\psi\\r\end{bmatrix}
 =
 \underbrace{\begin{bmatrix}\dL&\Lexc\\
                 \Lexc&-\dL\end{bmatrix}}_{\matr T_\psi}
 \begin{bmatrix}i_d\\i_e\end{bmatrix},
 \qquad
 \matr T_\psi\trans\matr T_\psi=(\dL^2+\Lexc^2)\matr I.
\end{equation}
After division by $\sqrt{\dL^2+\Lexc^2}$, the transformation is orthogonal:
it rotates or reflects the plane into a torque-effective direction and a
torque-invisible complementary direction.  The transformation is nonsingular
unless both $\dL$ and $L_e$ vanish.

Torque becomes $\Te=\kt\iq\sflux$ and contains no $r$.  A nonzero
constant-torque curve is a hyperbola in
the $(s_{\psi},i_q)$ plane and is independent of $r$; in the original 3D space
it is therefore a hyperbolic cylinder, possibly oblique relative to the current
axes.  At zero torque it degenerates into two planes:
\begin{equation}
 \iq=0\quad\text{or}\quad\Lexc\iexc+\dL\id=0.
\end{equation}

The same result follows from the symmetric quadratic form
\begin{equation}
 \Te=\state\trans\Hm\state,\qquad
 \Hm=\frac{\kt}{2}\begin{bmatrix}
 0&\dL&0\\ \dL&0&\Lexc\\0&\Lexc&0
 \end{bmatrix}.\label{eq:torque-matrix}
\end{equation}
Its eigenvalues are
\begin{equation}
 0,\qquad \pm\frac{\kt}{2}\sqrt{\dL^2+\Lexc^2}.
\end{equation}
The opposite signs produce saddle curvature; the zero eigenvalue extends each
2D hyperbola along a free direction.  Thus ``bilinear'' is the algebraic cause,
``saddle'' describes the scalar torque landscape, and ``hyperbolic cylinder''
classifies a nonzero torque level surface.

At fixed excitation, torque contours in the $dq$ plane resemble PSM contours
with an effective excitation flux $L_e\bar i_e$.  At fixed $i_q$, constant
torque is a straight line in $(i_d,i_e)$.  These slices explain why one 3D
surface can appear as different curve families in different planes.

\begin{figure}[tb]
\centering
\begin{tikzpicture}
\begin{axis}[publication 3d,width=.86\linewidth,height=7.3cm,
 xlabel={$i_d$},ylabel={$i_q$},zlabel={$i_e$},
 xmin=-2,xmax=2,ymin=-2,ymax=2,zmin=-2,zmax=2]
 \addplot3[surf,shader=interp,opacity=.48,draw=Purple!55,
  domain=.45:2,y domain=-1.5:1.5,samples=16,samples y=13]
  ({y},{.8/x},{(x+.45*y)/.9});
 \addplot3[surf,shader=interp,opacity=.48,draw=Purple!55,
  domain=-2:-.45,y domain=-1.5:1.5,samples=16,samples y=13]
  ({y},{.8/x},{(x+.45*y)/.9});
\end{axis}
\end{tikzpicture}
\caption{A nonzero torque level is hyperbolic in the effective-flux and
$i_q$ coordinates and extends along the independent third direction.}
\label{fig:torque-surface}
\end{figure}


\paragraph{Synthesis.}
The coordinate $s_\psi$ is selected by the physics already present in the
torque equation; $r$ is selected by orthogonality.  Algebra turns torque into
a product, geometry turns its nonzero levels into hyperbolas extended along a
free direction, and control design learns which current combination changes
torque and which can be traded without changing it.
